\section{Centered H\'enon warps and the exact classical clock}
\label{sec:model-clock}

The construction begins with a radial potential whose classical phase volume
can be evaluated exactly.  A determinant-one configuration map then changes
the force field without changing configuration volume.  This section isolates
that elementary invariant before any quantum or statistical interpretation is
introduced.

\subsection{The operator family}

For \(a>-1\), define the positive fixed point
\begin{equation}
 r_a=\frac{1}{1+\sqrt{1+a}}
\end{equation}
of the area-preserving H\'enon map
\begin{equation}
 H_a(x,y)=(1-ax^2-y,x).
\end{equation}
Translation by \((r_a,r_a)\) gives the centered conjugate
\begin{equation}
 \widetilde H_a(x,y)=(-c_ax-ax^2-y,x),
 \qquad c_a=2ar_a=2(\sqrt{1+a}-1).
 \label{eq:centered-henon}
\end{equation}
Its inverse and Jacobian are
\begin{equation}
 \widetilde H_a^{-1}(u,v)=(v,-c_av-av^2-u),
 \qquad
 D\widetilde H_a(x,y)=
 \begin{pmatrix}-c_a-2ax&-1\\1&0\end{pmatrix},
\end{equation}
so \(\det D\widetilde H_a=1\).  Both the map and its inverse are
polynomial.  Every fixed iterate
\(\Psi_{a,n}=\widetilde H_a^n\) is therefore a proper, area-preserving
polynomial automorphism.  When \(a\ne0\), its degree and the degree of its
inverse are \(D=2^n\).

Let
\begin{equation}
 V_{a,n}(q)=2\pi\exp\!\left(\pi|\Psi_{a,n}(q)|^2\right),
 \qquad
 A_B(x,y)=\frac{B}{2}(-y,x),
 \label{eq:potential-vector}
\end{equation}
and consider the Friedrichs realization on \(L^2(\R^2)\) of
\begin{equation}
 \boxed{
 \HH_{a,n,B}
 =\frac12(-\ii\nabla-A_B)^2+V_{a,n}.}
 \label{eq:main-operator}
\end{equation}
The parameter \(B\in\R\) is fixed.  Three values of \(a\) have distinct
roles.  The limit \(a=0\) gives
\(\widetilde H_0(x,y)=(-y,x)\), hence the exactly radial integrable control
\(V_{0,n}(q)=2\pi e^{\pi|q|^2}\).  The value \(a=1.02\) was fixed by an
earlier H\'enon study before the present operator was proposed; it is the
primary nonlinear test case, not a parameter selected by the theorem below.
The value \(a=6\) is retained as a high-distortion and independently motivated
hyperbolic control.

\subsection{A general clock-invariance theorem}

The H\'enon formula is not needed for the classical identity.  The invariant
is the pushforward of configuration area.

\begin{theorem}[Exact clock under area-preserving warps]
\label{thm:classical-clock}
Let \(\Psi:\R^2\to\R^2\) be a proper \(C^1\) diffeomorphism with
\(|\det D\Psi|=1\), and set
\(V_\Psi(q)=2\pi e^{\pi|\Psi(q)|^2}\).  For any vector potential \(A(q)\),
assumed measurable,
the classical symbol
\begin{equation}
 h_{\Psi,A}(q,p)=\frac12|p-A(q)|^2+V_\Psi(q)
\end{equation}
has normalized phase volume
\begin{equation}
 \Ncl(E)
 :=\frac{1}{(2\pi)^2}
 \operatorname{vol}\{(q,p):h_{\Psi,A}(q,p)<E\}
 =\rvmlaw+1
 \label{eq:exact-classical-clock}
\end{equation}
for \(E\ge2\pi\).  For \(E<2\pi\), the phase volume is zero.
\end{theorem}

\begin{proof}
At fixed \(q\), the momentum translation \(p'=p-A(q)\) preserves Lebesgue
measure.  The momentum disk has area \(2\pi(E-V_\Psi(q))_+\), so
\begin{equation}
 \Ncl(E)=\frac{1}{2\pi}\int_{\R^2}(E-V_\Psi(q))_+\dd q.
 \label{eq:momentum-reduction}
\end{equation}
The change \(u=\Psi(q)\) preserves configuration measure.  Writing
\(L(E)=\log(E/2\pi)\) and \(R(E)=\sqrt{L(E)/\pi}\), polar integration gives
\begin{align}
 \Ncl(E)
 &=\frac{1}{2\pi}\int_{|u|<R(E)}
       \left(E-2\pi e^{\pi|u|^2}\right)\dd u \\
 &=\int_0^{R(E)}\left(E-2\pi e^{\pi r^2}\right)r\dd r
 =\frac{E}{2\pi}L(E)-\frac{E}{2\pi}+1.
\end{align}
\end{proof}

Equivalently, if \(\Phi_\Psi(q)=\pi|\Psi(q)|^2\), then for \(t\geq0\),
\begin{equation}
 A_{\Phi_\Psi}(t)
 :=|\{q:\Phi_\Psi(q)<t\}|=t.
 \label{eq:sublevel-invariant}
\end{equation}
For \(t<0\), the sublevel area is zero.
The exact identity is independent of the geometry of individual level sets.
It also does not require the momentum translation to be a global canonical
transformation.  For a magnetic vector potential with nonzero curl, it
generally is not; only the fiberwise measure statement is used.

\begin{remark}[What the constant means]
The \(+1\) in \cref{eq:exact-classical-clock} belongs to the exact classical
comparison function.  The quantum remainder proved in \cref{sec:quantum-weyl}
is much larger than a constant.  No quantum \(7/8\) term, and no analogue of
the oscillatory \(S(T)\) term, is inferred from this identity.
\end{remark}

\Cref{fig:equal-area-geometry} displays the preimages of two radial circles.
The domains differ sharply in shape while their area remains exactly
\(L(E)\), as required by \cref{eq:sublevel-invariant}.

\begin{figure*}[t]
  \centering
  \includegraphics[width=0.96\textwidth]{fig2_equal_area_geometry.pdf}
  \caption{Configuration boundaries for two common energy levels.  The radial,
  centered \(a=1.02\), and high-distortion \(a=6\) domains have different
  geometry but exactly the same area at each energy by the determinant-one
  H\'enon change of variables.}
  \label{fig:equal-area-geometry}
\end{figure*}
